%% Chapter 2: Computational Methods
%% FINAL VERSION (surgical voice pass)
%%
%% UNRESOLVED CONFLICTS (author must confirm before submission):
%%   [CONFLICT-T]   Temperature: chapter uses 300 K;
%%                  LAMMPS scripts consistently use 298 K.
%%   [CONFLICT-TD]  Langevin damping: chapter uses tau=1,000 fs;
%%                  LAMMPS scripts use 100,000 fs (= 100 ps).
%%   [CONFLICT-SA]  SA schedule: see TODO in Section 2.3.4.
%%   [CONFLICT-RL]  Run length: chapter states 1,000 ns;
%%                  trajectory data indicates ~3.6 µs.

%% CHANGE 1 (opening paragraph) — your verbatim wording
Modeling ASO binding to mRNA required both a physically stable representation of the target RNA and a simulation framework where ASOs interact a target.
In this chapter, we describe the modeling pipeline. Experimental 3-dimensional structures or NuFold~\cite{Kagaya2025} predicted RNA structures were prepared; coarse-grained molecular dynamics (CGMD) was then used to fold and equilibrate each structure to determine the regions most likely to bind ASOs by assessing their secondary structure and Solvent-Accessible Surface Area (SASA). Finally, ASO binding was quantified based on centre of mass (COM) distance. All simulations were run with LAMMPS as before, using the SIS coarse-grained RNA force field~\cite{Thompson2022} with an implementation from Aierken and Joseph~\cite{Aierken2024} for RNA folding and binding studies. Subsequently, post-processing was carried out with custom Python codes and the tools described above.
%% - - - - - - - - - - - - - - - - - - - - - - - - - - - - - - - - - - -
\section{RNA Sequences and System Preparation}
\label{sec:system_prep}
%% - - - - - - - - - - - - - - - - - - - - - - - - - - - - - - - - - - -

\subsection{3NJ6 Reference Structure}
\label{sec:3nj6}

Validation of our simulation parameters was carried out using an existing crystal structure of a CAG-repeat RNA duplex 3NJ6~\cite{Aierken2024}. Because this structure has already been validated against SIS model’s parameters, we used it as a benchmark to test whether our force field has been properly implemented. To create the validation structure, a single strand was used, with the extract of C3$^\prime$ atom coordinates serving as a fixed reference for calculating root mean square deviations (RMSD).

\subsection{NuFold Benchmark Structures}
\label{sec:nufold}

In addition to the 3NJ6 benchmark, we tested the sequences from the results of Figure~5 of Kagaya et al.~\cite{Kagaya2025}, which evaluate how well NuFold predicts structures of a variety of RNA types. The seven sequences that were used were 2DER (C chain, 76 nts); 6CK4 (A chain, 117 nts); 4QLN (rp12, 125 nts); 5ML7 (A chain, 96 nts); 4OJI (A chain, 54 nts); 6G7Z (rp05); and 4U7U (L chain, 61 nts). 
In order to predict the secondary structure, we first performed multiple sequence alighnement using rMRSA on the target's FASTA, and the used IPknot to perform secondary structure constraing predictions, and thirdly, Nufold is executed using the training configeration preset. The rank-1 output PDB for each sequence represented the predicted three dimensional structure for the corresponding RNA. Each rank-1 PDB was subsequently read in order to obtain its C3$^\prime$ coordinates which would be stored as the fixed RMSD reference for all subsequent simulation runs for that specific RNA target.


\subsection{Longer and mRNA-Scale Sequences}
\label{sec:longer_seqs}

Ten RNA sequences were randomly generated and sampled (50-200 bases). Additionally, sequences of approximately 1000 bases were used to test the ability of the pipeline to scale to mRNA-length chains. All sequence within this length range were folded according to the Simulated Annealing protocol outlined in Section~\ref{sec:sa}, and secondary structures were analyzed with RNAfold as explained in Section~\ref{sec:sec_struct}.

\subsection{Coarse-Grained Topology Generation}
\label{sec:topology}

\subsubsection{Sequence Extraction from PDB Files}

For each RNA target, the source PDB or PDB1 file was parsed to extract
RNA residue information.
ATOM and HETATM records were filtered to retain only recognized RNA nucleotide residue.
When multiple chains were present in a PDB file, the chain containing
the greatest number of RNA residues was selected automatically; for
the NuFold targets, chain selections were fixed to match the
assignments in Kagaya et al.~\cite{Kagaya2025}
(e.g., chain~C for 2DER, chain~A for 6CK4).
%% CHANGE 3 ("for downstream use") — audit fix: name the actual next step
Residues were sorted in ascending sequence-number order; the resulting
nucleotide sequence (A, C, G, U) was used as the basis for LAMMPS data
file generation.


The PDB files for all RNA molecules were analyzed to obtain inforamtion about the RNA residues. After sorting the residues in increasing order of their sequence numbers, the resulting nucleotide sequence (A, C, G, U) was used to generate the LAMMPS data files.


\subsubsection{LAMMPS Data File Generation}

The initial LAMMPS configeration places the beads leanearly along the $x$-axis at a with a 5.9~\AA \space spacing between beads -- equal to the harmonic bond equilibrium length. A fully extended chain is produced irrespective of the geometry of the source PDB file. Sequential 1-2 pair bonds and sequential 1-2-3 triplet angles are created according to the residue order. Three dimensional periodic boundary conditions are applied to the simulation box that is a cube with variable input edge lengths. Types of atoms, masses of atoms, positions of atoms at start time, connections of atoms via bonds, connections of atoms via angles, and lengths of edges of the simulation box are recorded into a LAMMPS configuration. 


The C3$^\prime$ atom from the source PDB files are obtained used as a reference structure to calculate RMSDs; these atomic coordinates are not used as initial bead locations within the LAMMPS configuration.


%% - - - - - - - - - - - - - - - - - - - - - - - - - - - - - - - - - - -
\section{Coarse-Grained Force Field}
\label{sec:forcefield}
%% - - - - - - - - - - - - - - - - - - - - - - - - - - - - - - - - - - -

All simulation methods used the SIS coarse-grained RNA force field developed by Nguyen, Hori, and Thirumalai~\cite{Aierken2024}, with a simulation pipeline developed by Aierken and Joseph~\cite{Aierken2024}. Each nucleotide was represented as a bead with interactions that have been parametrized to represent the relative energies associated with each type of standard Watson-Crick base pair interaction (C-G and A-U) and G-U wobble base pair interaction. The total potential energy is given by:

\begin{equation}
  U = U_{\mathrm{BA}} + U_{\mathrm{EV}} + U_{\mathrm{BP}},
  \label{eq:Utot}
\end{equation}

where $U_{\mathrm{BA}}$ represents the bonded backbone potentials,
$U_{\mathrm{EV}}$ is the excluded-volume potential due to repulsion, and $U_{\mathrm{BP}}$
is the many-body base-pairing potential.


The solvent is modeled implicitly; Therefore, all ASO -RNA simulations were conducted under an effective an effective ionic strength of 50~mM NaCl consistent with the
conditions used to benchmark the COM-based $K_D$ method using single-molecule fluorescence resonance energy transfer (FRET) data from the SARS-CoV-2 nucleocapsid system~\cite{Alston2023}.






\subsection{Bonded Interactions}

The backbone covalent connectivity was modeled with a harmonic bond potential,
\begin{equation}
  U_{\mathrm{bond}} = k(r - r_0)^2,
\end{equation}
with a spring constant of $k = 7.5$~kcal\,mol$^{-1}$\,\AA$^{-2}$ and
equilibrium length of $r_0 = 5.9$~\AA.
The local stiffness of the backbone is captured by a harmonic angle potential,
\begin{equation}
  U_{\mathrm{angle}} = k_\theta(\theta - \theta_0)^2,
\end{equation}
with bending constant $k_\theta = 5.0$~kcal\,mol$^{-1}$\,rad$^{-2}$
and equilibrium angle $\theta_0 = 150.0^\circ$, consistent with the
extended geometry of single-stranded nucleic acid backbones.

We also implement a standard Lennard-Jones 12 - 6 potential
with $\varepsilon = 2.0$~kcal\,mol$^{-1}$ and $\sigma = 8.9090$~\AA\
is shifted to zero at a cutoff of 10.0~\AA\ via the
\texttt{pair\_modify shift} command, retaining only the repulsive core.
The large $\sigma$ reflects the physical size of a nucleotide-level bead.

\subsection{Base-Pairing Potential}

The base pairing and stacking interactions in this system are modeled using a custom LAMMPS pair style \texttt{base/rna}~\cite{Aierken2024}, with a cutoff of 18.0~\AA. This pair style represents the anisotropic geometry of Watson-Crick base pairing using inter-bead distances, four angles, and two dihedral angles to reproduce the A-form helix of RNA. Interaction strengths are: C--G pairs at $\varepsilon = -5.0$~kcal\,mol$^{-1}$,
A--U and G--U wobble pairs at $\varepsilon = -3.33$~kcal\,mol$^{-1}$,
and all non-complementary pairs at $\varepsilon = 0$.
All types of base pairs have geometric parameters: the equilibrium contact distance of
3.0~\AA, outer cutoff distance of 13.8~\AA, and shape parameters of
$(1.5,\;0.5,\;1.8326,\;0.9425,\;1.8326,\;1.1345)$.





Non-bonded interactions between atoms that are separated by less than one bond length (i.e., 1-2 and 1-3) are completely excluded from being calculated in order to prevent double-counting of bonded interactions. However, all 1-4 non-bonded interactions are still included at their full strength (\texttt{special\_bonds lj 0.0 0.0 1.0}). Neighbor lists are created using a bin algorithm with a skin distance of 15.0~\AA\ and are updated every 10 timesteps.




%% - - - - - - - - - - - - - - - - - - - - - - - - - - - - - - - - - - -
\section{Simulation Protocol}
\label{sec:sim_protocol}
%% - - - - - - - - - - - - - - - - - - - - - - - - - - - - - - - - - - -

\subsection{Energy Minimization}

Each System undergoes Conjugate Gradient Energy Minimization before any Dynamics are conducted on it. The Force and Energy tolerance levels used in this process were $10^{-4}$ kcal mol$^{-1}$ \AA$^{-1}$ and $10^{-6}$ kcal mol$^{-1}$. Also, the force evaluation limit was set at 10,000. This is done to remove steric clashes generated by Topology Generation using Coarse Graining.

\subsection{\textit{NVT} Equilibration}
Each minimized system was run for 50~ns ($5 \times 10^{6}$ time steps and $\Delta t=10$ fs) in the canonical ensemble (\textit{NVT}), as described above. The ASO simulations were also performed at temperatures of $298$ K and $300$ K.

The \textit{NVT} ensemble uses the NVE velocity-verlet integration method, coupled with a Langevin thermostat~\cite{Schneider1978}. The thermal coupling constant used by the Langevin thermostat has been set to $\tau = 100,000$ fs.

The initial velocities were all sampled from a Gaussian distribution around the target temperature. The linear momentum associated with each system's center of mass was zeroed every 1,000 steps.

\subsection{Production Runs}
\label{sec:production}

Sampling was done using production runs for both RNA and ASO under \textit{NVT} conditions, but over different time scales: 1,000~ns for ASO simulations and 4~\textmu s for RNA simulations. 
All restart files were generated every $10^7$ steps, and the final configuration was written to a LAMMPS-formatted data file. 
Positions of all atoms were recorded every 50,000 timesteps (500~ps) with position information including ID’s for molecule/atom, types of molecules/atoms, charges on the atoms, and unwrapped cartesian coordinates (\texttt{xu yu zu}).
At least 5 independent production runs were performed for each system utilized in calculating binding affinities (section~\ref{sec:binding_affinity}).


\subsection{Enhanced Sampling via Simulated Annealing}
\label{sec:sa}

Standard \textit{NVT} molecular dynamic simulations for nucleic acid 
sequences having stable secondary structural elements are prone to kinetic 
trapping prior to sampling the entire set of possible conformations. 
%% CHANGE 5 (SA transition) — your exact wording; footnote adjusted to your wording
Simulated annealing (SA) has been employed on both the NuFold 
benchmark sequences and on the longer (200~nt or greater) random sequences.
\footnote{Although the NuFold benchmark sequence lengths (54-125~nt) do not meet 
the length criteria (200~nt) established for the random sequences, they have 
compact, well defined tertiary structures that provide a realistic opportunity 
for kinetic trapping during \textit{NVT} simulation. Thus, simulated annealing was utilized.} 

I carried out a 10 step heating cycle that consisted of: ramp 298~K$\rightarrow 500~K$ (2.5~ns) $\rightarrow$ hold at $500~K (2.5~ns) $\rightarrow$
cool to 280~K (5~ns) $\rightarrow$ hold at 280~K (2.5~ns);

The convergence of SA trajectories versus \textit{NVT} runs was evaluated by comparison of
convergence diagnostic values including $R_g$, RMSD, and potential energy.






%% - - - - - - - - - - - - - - - - - - - - - - - - - - - - - - - - - - -
\section{Trajectory Analysis}
\label{sec:traj_analysis}
%% - - - - - - - - - - - - - - - - - - - - - - - - - - - - - - - - - - -

\subsection{Thermodynamic Observables}


We collected values of potential energy, kinetic energy, total energy, temperature, pressure, and mass density at every 10,000 steps (100~ps) during the simulation. 
The time averages are collected through block averaging, where instantaneous values are sampled every 10,000 steps and then concatenated into blocks of 10, written to an output file every 100,000 steps.


\subsection{Structural Observables: $R_g$, RMSD, and RMSF}
\label{sec:structural_obs}

The radius of gyration $R_g$, backbone bead RMSD, and per-nucleotide
root-mean-square fluctuation (RMSF) are computed from saved trajectory
frames using the C3$^\prime$-based reference coordinates described in
Sections~\ref{sec:3nj6} and~\ref{sec:nufold}.

$R_g$ serves as the primary order parameter for monitoring global
conformational state: a decreasing $R_g$ relative to the initial
extended chain indicates compaction associated with secondary structure
formation.

%% CHANGE 6 (Kabsch) — your verbatim wording
RMSD was computed at each saved frame after translating both the
trajectory frame and the reference structure to their respective centers
of mass and applying the Kabsch algorithm to obtain the optimal
rigid-body superposition.
For the 3NJ6 validation system, the average production-run RMSD is
compared against the value reported by Aierken and
Joseph~\cite{Aierken2024} to confirm correct force-field implementation.

RMSF is computed per nucleotide bead over the aligned production
trajectory.
For each bead $i$,

\begin{equation}
  \mathrm{RMSF}_i = \sqrt{\frac{1}{F}\sum_{t=1}^{F}
    \left|\mathbf{r}_i(t) - \langle\mathbf{r}_i\rangle\right|^2},
  \label{eq:rmsf}
\end{equation}

where $\mathbf{r}_i(t)$ is the Kabsch-aligned position of bead $i$ at
frame $t$, $\langle\mathbf{r}_i\rangle$ is its time-averaged position,
and $F$ is the total number of sampled production frames.
%% CHANGE 7 (RMSF downstream use) — your verbatim wording
RMSF profiles were plotted by residue index; residues with elevated RMSF
were treated as conformationally flexible and compared against the
secondary-structure and SASA analyses when interpreting site
accessibility.



\subsection{Thermodynamics Observable }

\subsection{Structural Observables: Rg, RMSD & RMSF} \label{sec:structural_obs}
Radius of Gyration Rg , Backbone Bead RMSD & Per Nucleotide Root Mean Square Fluctuation (RMSF), are calculated from Trajectory Frames Saved Using Reference Coordinates C3’ Based (see Section~\ref{sec:3nj6} & Section~\ref{sec:nufold}) .
The Radius of Gyration Rg will be the Primary Order Parameter for Global Conformational State. A Decrease in Radius of Gyration Rg Compared to Initial Extended Chain Indicates Compaction due to Secondary Structure Formation.
%% CHANGE 6 (Kabsch) — your verbatim wordingRMSD was computed at each saved frame after translating both the trajectory frame and the reference structure to their respective centers of mass and applying the Kabsch algorithm to obtain the optimal rigid-body superposition. For the 3NJ6 validation system, the average production run RMSD will be compared to the value of the Aierken and Joseph report~\cite{Aierken2024} to confirm a Correct Implementation of Force Field.
RMSF is computed per nucleotide bead along the aligned production trajectory. For each bead i:
\begin{equation}\mathrm{RMSF}_i = \sqrt{\frac{1}{F}\sum_{t=1}^{F}\left|\mathbf{r}_i(t) - \langle\mathbf{r}_i\rangle\right|^2},\label{eq:rmsf}\end{equation}
where r_i(t) is the Position of Kabsch Aligned Bead i at Frame t, \langle r_i \rangle is it's time averaged position, and F is the Total Number of Sampled Frames in the Production Run.






%% - - - - - - - - - - - - - - - - - - - - - - - - - - - - - - - - - - -
\section{Secondary Structure Analysis}
\label{sec:sec_struct}
%% - - - - - - - - - - - - - - - - - - - - - - - - - - - - - - - - - - -

\subsection{Minimum Free Energy Prediction with ViennaRNA}

Secondary structure predictions were computed with the ViennaRNA
Package 2.0~\cite{Lorenz2011}.
\texttt{RNAfold} was used to generate the minimum free energy (MFE)
structure and base-pair probability matrix for each sequence using
McCaskill's partition function algorithm with Turner 2004
nearest-neighbor parameters.
For sequences longer than 400~nt, \texttt{RNALfold} was used instead
to compute locally stable substructures within a restricted base-pair
span, making secondary structure analysis tractable at mRNA scale.
Dot-bracket output from both programs was parsed to extract paired and
unpaired nucleotide positions across all sequences.

\subsection{Identification of Single-Stranded Open Sites}

Nucleotide positions classified as unpaired in the MFE structure are
designated \emph{open sites} and treated as candidate regions for ASO
hybridization.
Open sites were identified from ViennaRNA dot-bracket output and
tabulated by sequence position and contiguous run length.
These positions define the primary 2D accessibility criterion and are
subsequently intersected with the 3D SASA criterion
(Section~\ref{sec:sasa}) to generate a ranked candidate set.

\subsection{Structure Visualization}

Two-dimensional secondary structure representations were generated using
the forna web tool, which renders RNA structure as a force-directed graph
with nucleotides as nodes and base pairs as edges.
forna diagrams were used to inspect folding topology, verify base-pair
assignments from ViennaRNA output, and produce secondary structure
figures.

%% - - - - - - - - - - - - - - - - - - - - - - - - - - - - - - - - - - -
\section{Solvent-Accessible Surface Area Analysis}
\label{sec:sasa}
%% - - - - - - - - - - - - - - - - - - - - - - - - - - - - - - - - - - -

Secondary structure analysis identifies nucleotides that are unpaired in
two dimensions, but a nucleotide can be unpaired yet sterically buried
in a compact tertiary fold.
SASA calculations address this gap: per-nucleotide SASA was computed
from a representative set of evenly spaced production-run frames by
rolling a probe sphere over the molecular surface.
\textbf{[ADD: specific SASA tool FreeSASA]}
The resulting per-nucleotide values were time-averaged and plotted as a
function of sequence position.
Nucleotides with above-average SASA are treated as three-dimensionally
accessible; those below the average are considered buried and excluded
from the high-confidence candidate set.

%% - - - - - - - - - - - - - - - - - - - - - - - - - - - - - - - - - - -
\section{Integrated Binding Site Identification}
\label{sec:site_id}
%% - - - - - - - - - - - - - - - - - - - - - - - - - - - - - - - - - - -

Candidate ASO target sites are identified by the intersection of the 2D
and 3D accessibility criteria: a site is retained as high-confidence
only when it is unpaired in the MFE structure \emph{and} has
above-average per-nucleotide SASA across the production run.
Sites that are open in 2D but buried in 3D are deprioritized as
geometrically inaccessible to an incoming ASO molecule.
The surviving candidates are passed to the binding affinity calculations
described in Section~\ref{sec:binding_affinity}.

%% - - - - - - - - - - - - - - - - - - - - - - - - - - - - - - - - - - -
\section{Quantifying ASO -RNA Binding Affinity}
\label{sec:binding_affinity}
%% - - - - - - - - - - - - - - - - - - - - - - - - - - - - - - - - - - -

%% CHANGE 8 (binding affinity section opening) — your verbatim wording
After candidate binding regions had been identified, ASO -RNA
interactions were quantified directly from CGMD trajectories as
apparent dissociation constants ($K_D$).
These values were computed using the COM distance framework of Alston,
Soranno, and Holehouse~\cite{Alston2023}, as validated against the
$K_D$ formulation of Tesei et al.~\cite{Tesei2021}.

\subsection{Classification of Bound and Unbound States}
\label{sec:bound_unbound}

%% CHANGE 9 (minimum-image) — audit fix approved by you
For each trajectory, the minimum distance between any ASO bead and any
RNA bead is computed at every saved frame.
Because trajectories are dumped as unwrapped Cartesian coordinates
(\texttt{xu yu zu}), minimum-image correction is applied to recover
the true periodic distance before classification.
This distance distribution typically follows a bimodal profile: a
short-distance peak corresponding to the bound state and a long-distance
peak corresponding to the free, diffuse state~\cite{Alston2023}.
Gaussian functions are fit to both peaks, and the threshold separating
bound from unbound is placed at the distance that minimizes the overlap
between the two distributions.
Because the optimal threshold depends on the physical dimensions of
both molecules, it is determined independently for each ASO -RNA pair.
A frame is classified as bound only when five or more consecutive frames
satisfy the distance criterion, filtering out brief stochastic collisions
that do not represent stable hybridization events.

\subsection{Binding Persistence and Closest-Approach Analysis}

%% CHANGE 10 (binding persistence opener) — audit fix: drop meta-announcement
\emph{Binding persistence} is defined as the number of production-run
frames classified as bound for each free ASO molecule  -- a direct count
of contact time.

For systems where binding is sparse and $f_\text{bound}$ is
near-quantized, a \emph{closest-approach cumulative distribution
function} (CDF) provides a graded affinity ranking.
For each free ASO molecule $k$, the minimum RNA -ASO distance achieved
over the entire trajectory, $d_{\min,k}$, is recorded.
The CDF

\begin{equation}
  P(d_{\min} \leq d) = \frac{1}{N_\text{ASO}}
    \sum_{k=1}^{N_\text{ASO}} \mathbf{1}[d_{\min,k} \leq d]
  \label{eq:cdf}
\end{equation}

%% CHANGE 11 (CDF closing) — audit fix: split into two clean sentences
is plotted on a log scale as a function of the distance threshold $d$.
This provides a graded affinity ranking based on closest-approach
geometry, which is more informative than binding persistence alone when
stable binding events are rare.

\subsection{Apparent Dissociation Constant}
\label{sec:kd}

Let $f_{\mathrm{bound}}$ denote the fraction of production-run frames
classified as bound.
The apparent $K_D$ is then

\begin{equation}
  K_D = \frac{(1 - f_{\mathrm{bound}})^2}{N_A \, f_{\mathrm{bound}} \, V},
  \label{eq:KD}
\end{equation}

where $N_A$ is Avogadro's number and $V$ is the simulation box volume
in liters, giving $K_D$ in mol\,L$^{-1}$ (M).
This expression follows from the law of mass action for a two-state
binding equilibrium~\cite{Tesei2021}.
Because the SIS force field encodes effective free energies at
coarse-grained resolution, $K_D$ values are used comparatively to rank
ASO variants rather than as absolute predictions of solution-phase
affinity.

\subsection{Normalized Binding Affinity ($K_A^*$)}
\label{sec:KA_star}

All $K_D$ values are normalized by the $K_D$ of the unmodified
reference system (\texttt{25ASO\_unmod} binding to the target
single-chain RNA).
The normalized association constant is defined as

\begin{equation}
  K_A^* = \frac{K_A(\text{variant})}{K_A(\texttt{25ASO\_unmod})},
  \qquad K_A = 1/K_D.
  \label{eq:KA_star}
\end{equation}

$K_A^* > 1$ indicates tighter binding than the unmodified reference;
$K_A^* < 1$ indicates weaker binding.
This normalization follows Alston et al.~\cite{Alston2023}, in which
simulated $K_A^*$ values showed semi-quantitative agreement with
single-molecule FRET experiments when both were expressed relative to
an internal reference.

Error propagation for $K_A^*$ follows standard ratio error propagation:

\begin{equation}
  \sigma(K_A^*) = K_A^* \sqrt{\left(\frac{\sigma_A}{A}\right)^2
    + \left(\frac{\sigma_B}{B}\right)^2},
  \label{eq:error_prop}
\end{equation}

where $A$ and $B$ are the $K_A$ values of the variant and reference,
and $\sigma_A$, $\sigma_B$ are their standard errors across a minimum
of five independent simulation replicates.

